% !TEX program = pdflatex
% !TEX encoding = UTF-8 Unicode

% Plantilla, basada en la clase `scrbook` del paquete KOMA-script,  para la elaboración de un TFG siguiendo las directrices del la comisión del Grado en Matemáticas de la Universidad de Granada.

% Francisco Torralbo Torralbo

\documentclass[print, color]{ugrTFG}

% VERSIÓN ELECTRÓNICA PARA TABLETA
% Cambiando la opción "print" por "tablet" generaremos un pdf adaptado para leerlo en tabletas de 9 pulgadas.

% -------------------------------------------------------------------
% INFORMACIÓN DEL TFG Y EL AUTOR
% -------------------------------------------------------------------

\newcommand{\miTitulo}{Combinando distintas técnicas para el diseño de una metahehurística para problemas de optimización de alta densidad \xspace}
\newcommand{\miNombre}{Jesús García León\xspace}
\newcommand{\miGrado}{Doble Grado en Ingeniería Informáticva y Matemáticas}
\newcommand{\miFacultad}{Escuela Técnica Superior de Ingeniería Informatica y de Telecomunicación y Facultad de Ciencias}
\newcommand{\miUniversidad}{Universidad de Granada}

% Añadir tantos tutores como sea necesario separando cada uno de ellos mediante el comando `\medskip` y una línea en blanco
\newcommand{\miTutor}{Daniel Molina Cabrera \\ \emph{Departamento de Ciencias de la Computación e Inteligencia Artificial} 

  % Añadir tantos tutores como sea necesario. 

  %\medskip
  %Nombre del tutor 2 \\ \emph{Departamento del tutor 2}
}
\newcommand{\miCurso}{2023-2024\xspace}

\hypersetup{
	pdftitle={\miTitulo},
	pdfauthor={\textcopyright\ \miNombre, \miFacultad, \miUniversidad}
}

\begin{document}

\maketitle

% -------------------------------------------------------------------
% FRONTMATTER
% -------------------------------------------------------------------
\frontmatter % Desactiva la numeración de capítulos y usa numeración romana para las páginas

% !TeX root = ../tfg.tex
% !TeX encoding = utf8
%
%*******************************************************
% Declaración de originalidad
%*******************************************************

\thispagestyle{empty}

\hfill\vfill

\textsc{Declaración de originalidad}\\\bigskip

D./Dña. \miNombre \\\medskip

Declaro explícitamente que el trabajo presentado como Trabajo de Fin de Grado (TFG), correspondiente al curso académico \miCurso, es original, entendido esto en el sentido de que no he utilizado para la elaboración del trabajo fuentes sin citarlas debidamente.
\medskip

En Granada a \today 
\vspace{3cm}
\begin{center} 
Fdo: \miNombre 

\end{center}

\vfill

\cleardoublepage
\endinput
   
% !TeX root = ../tfg.tex
% !TeX encoding = utf8

%*******************************************************
% Dedication
%*******************************************************
\thispagestyle{empty}
\phantomsection 
\pdfbookmark[1]{Dedicatoria}{Dedicatoria}

\hfill
\vfill

\begin{flushright}
\itshape
Dedicatoria (opcional) \\
Ver archivo \texttt{preliminares/dedicatoria.tex}
\end{flushright}

\vfill

\cleardoublepage
\endinput
                % Opcional
% !TeX root = ../tfg.tex
% !TeX encoding = utf8

%*******************************************************
% Table of Contents
%*******************************************************
\phantomsection
\pdfbookmark[0]{\contentsname}{toc}

\setcounter{tocdepth}{2} % <-- 2 includes up to subsections in the ToC
\setcounter{secnumdepth}{3} % <-- 3 numbers up to subsubsections

\tableofcontents 

%*******************************************************
% List of Figures and of the Tables
%*******************************************************

    % *******************************************************
    %  List of Figures
    % *******************************************************    
    \phantomsection 
    % \listoffigures

    %*******************************************************
    % List of Tables
    %*******************************************************
    \phantomsection 
    % \listoftables
    
    %*******************************************************
    % List of Listings
    % The package \usepackage{listings} is needed
    %*******************************************************      
	  % \phantomsection 
    % \renewcommand{\lstlistlistingname}{Listados de código}
    % \lstlistoflistings 

\cleardoublepage
            
% !TeX root = ../tfg.tex
% !TeX encoding = utf8

%*******************************************************
% Agradecimientos
%*******************************************************

\chapter{Agradecimientos}

Agradecimientos (opcional, ver archivo \texttt{preliminares/agradecimiento.tex}).

\cleardoublepage
\endinput
            % Opcional

% !TeX root = ../tfg.tex
% !TeX encoding = utf8
%
%*******************************************************
% Summary
%*******************************************************

\selectlanguage{english}
\chapter{Abstract}

\indent The increase in computational power of computers has enabled the resolution of increasingly complex problems that depend on a growing number of variables. The use of classical algorithms in this context is either infeasible or extremely costly due to the exponential growth in difficulty brought by high dimensionality. This makes it necessary to design algorithms capable of finding solutions to high-dimensional problems. This issue is known as the \textit{curse of high dimensionality}. This phenomenon refers to the exponential growth of the search space as the dimensionality of a problem increases, which severely hinders the ability to find solutions.

In high-dimensional problems, many techniques that work efficiently in low-dimensional spaces become ineffective, as data occupies only a small fraction of the total space. This impacts search, classification, and optimization algorithms, as the growth in the search space leads to an increase in the required computation time, making it impractical to evaluate all possible options.

In this work, we propose the study and implementation of various \textbf{variable grouping} techniques, which aim to separate the variables of a problem into independent subgroups, allowing each subgroup to be optimized separately. This significantly reduces the effective dimensionality of the problem, enabling more efficient optimization. By dividing the set of variables into independent subgroups, we leverage parallelism and reduce computational complexity, as each subgroup can be optimized without affecting the others. This technique is particularly useful in problems where certain variables are highly correlated with each other but have little or no relationship with other variables in the set.

We will study this technique by applying it to algorithms designed to solve high-dimensional problems as well as algorithms that perform better in low dimensions. This will allow us to compare whether it is more efficient to design algorithms that improve performance in high dimensions or to divide the problem into smaller subproblems. To this end, we will introduce the mathematical concepts necessary to understand how to solve minimization problems, as well as the theory behind variable grouping algorithms. Additionally, we will analyze which statistical tests can be used to achieve an objective and effective comparison, explaining the rationale for their selection.

We will also present basic metaheuristic algorithms that we will later combine to develop our proposal. Among these are \textit{SHADE} and \textit{SHADE-ILS} as optimization algorithms, and \textit{ERDG} as a variable grouping algorithm.

Moreover, we will develop a library in the \textit{Julia} programming language to implement the algorithms used. This library aims not only to evaluate the effectiveness of the variable grouping technique but also to contribute to the \textit{Julia} community by providing a tool that other developers can use and improve in the future. To compare results, we will use the \textit{CEC2013 LSGO} benchmark, specifically designed for global optimization problems in high dimensions.

\textbf{Keywords:} computational power, optimization, neural networks, high dimensionality, curse of dimensionality, variable grouping, computational complexity, metaheuristics, Julia.


% Al finalizar el resumen en inglés, volvemos a seleccionar el idioma español para el documento
\selectlanguage{spanish} 
\endinput
                    
% !TeX root = ../tfg.tex
% !TeX encoding = utf8
%
%*******************************************************
% Introducción
%*******************************************************

% \manualmark
% \markboth{\textsc{Introducción}}{\textsc{Introducción}} 

\chapter{Introducción}

\indent Cuando enfrentamos un problema de optimización que depende de pocos factores, es sencillo idear estrategias para resolverlo y encontrar una solución. Sin embargo, a medida que crece el número de variables, la complejidad del problema aumenta, y se hace inviable resolverlo mediante métodos tradicionales. Para abordar esta problemática, se han diseñado algoritmos específicos para manejar la alta dimensionalidad utilizando metaheurísticas. Estas buscan optimizar la exploración del espacio de búsqueda, logrando un equilibrio entre exploración y explotación, lo que permite realizar búsquedas efectivas en un espacio donde la porción representada por la población de búsqueda es ínfima comparada con el espacio total. No obstante, estos algoritmos presentan limitaciones, y cuando el número de variables es suficientemente alto, hay dos opciones, por un lado, las metaheurísticas especialmente diseñado para LSGO, como SHADE-ILS, y por el otro, técnicas que permitan reducir la dimensionalidad que aborda el algoritmo.

En este trabajo, proponemos utilizar técnicas de \textbf{agrupamiento diferencial automático de variables}, que permite identificar las variables que interactúan entre sí, para asignarles el mismo grupo. De esta forma somos capaces de reducir el problema en subproblemas, uno para cada subgrupo ya que cada variable solo interactúa con las de su propio grupo y no (o mínimamente) con las de otros grupos. Así, podemos reducir la dimensión efectiva del problema, convirtiendo un problema de optimización de gran escala en problemas de optimización de menor tamaño. Nuestro objetivo es hibridar esta técnica con algoritmos diseñados para la alta dimensionalidad y con algoritmos que no están específicamente adaptados para ello. Compararemos el rendimiento aplicando y no aplicando esta técnica en ambos tipos de algoritmos.

Para probar la efectividad de esta técnica, se han planteado los siguientes objetivos:

\section*{Objetivos}

\subsection*{Parte matemática}
\begin{itemize}
    \item \textbf{Análisis teórico}: Repaso de las técnicas matemáticas tradicionales en la literatura para resolver problemas de optimización sin restricciones.
    \item \textbf{Estudio de teoremas}: Exposición y análisis de los teoremas en los que se basa el agrupamiento diferencial.
    \item \textbf{Fundamentos estadísticos}: Explicación de los tests estadísticos utilizados para la comparación de algoritmos.
\end{itemize}

\subsection*{Parte informática}
\begin{itemize}
    \item \textbf{Desarrollo de una biblioteca de técnicas de agrupamiento}: Implementación de algoritmos para optimización continua sin restricciones en el lenguaje de programación Julia y de técnicas de agrupamiento automático de variables.
    \item \textbf{Selección de algoritmos}: Descripción de los algoritmos que se incluirán en la biblioteca y que se emplearán en las comparaciones.
    \item \textbf{Evaluación experimental}: Análisis de los algoritmos en el benchmark \textit{cec2013lsgo}, presentación de resultados y conclusiones obtenidas.
\end{itemize}

\section*{Estructura de la memoria}

El trabajo se divide en tres partes claramente diferenciadas y cada parte está dividida en capítulos. Ademas de las dos partes que constituyen el contenido de este TFG, tenemos esta primera sección introductoria, donde se pone en contexto el trabajo, se definen los objetivos y se hace un repaso de la bibliografía necesaria. 
Además, se hace una estimación del presupuesto y del tiempo dedicado a cada tarea.

La parte Matemática está dividida en 3 capítulos, un primer capítulo donde se realizan definiciones importantes sobre optimización de funciones, se exponen algunos teoremas necesarios y se explican los métodos de búsqueda lineal y el algoritmo L-BFGS-B, que será un componente importante de nuestro algoritmo final. El segundo capítulo de esta sección se define el agrupamiento diferencial y se citan y demuestran los teoremas más importantes que justifican el diseño y utilidad de los algoritmos de agrupamiento automático de variables. Además se proporcionan algunas definiciones previas necesarias sobre separabilidad de funciones. Por último en la tercera sección de esta parte, se explican los fundamentos de los tests estadísticos y se definen los tests que suelen ser utilizados para comparar algoritmos. 

En la parte informática, tenemos un capitulo inicial donde se definen los componentes de los algoritmos que formarán nuestra propuesta, una segunda parte donde se explican los algoritmos que luego hibridaremos con el agrupamiento de variables. En tercer lugar se explicará nuestra propuesta: combinar ERDG con SHADE y SHADE-ILS. Por último, unas secciones finales donde se exponen los resultados obtenidos y las conclusiones que hemos obtenido de los resultados.

\chapter{Repaso Bibliográfico}

En esta sección, haremos un repaso bibliográfico sobre la optimización global a gran escala (LSGO, por sus siglas en inglés, \textit{Large Scale Global Optimization}). Se proporcionará un contexto general de las distintas técnicas metaheurísticas y enfoques que existen para abordar este tipo de problemas, y se situarán nuestros algoritmos en este contexto, explicando en mayor profundidad aquellas técnicas que están más relacionadas con los algoritmos \textit{SHADE} \cite{TanabeShade}, \textit{SHADE-ILS} \citep{Molina2018}, \textit{ERDG} \cite{ERDG} y \textit{DG2} \cite{DG2}. Para ello nos basaremos en el estudio realizado por \citeauthor{Review_I} en \cite{Review_I} y \cite{Review_II}.

\subsection{Enfoques Principales en LSGO}

A continuación, se presentan los enfoques principales en LSGO, los cuales se han desarrollado para afrontar la complejidad de los problemas de optimización de gran escala:

\begin{enumerate}
    \item \textbf{Descomposición del problema}: Dividir el problema en subproblemas manejables.
    \item \textbf{Hibridación y búsqueda local memética}: Combina algoritmos evolutivos con técnicas de búsqueda local.
    \item \textbf{Operadores de muestreo y variación}: Técnicas de muestreo para explorar el espacio de búsqueda.
    \item \textbf{Modelado por aproximación y uso de modelos sustitutos}: Se utilizan modelos simplificados para reducir el coste computacional.
    \item \textbf{Métodos de inicialización}: Métodos para asegurar una cobertura uniforme del espacio de búsqueda.
    \item \textbf{Paralelización}: Uso de múltiples instancias de algoritmos para acelerar la búsqueda.
\end{enumerate}

\subsection{Explotación de la Estructura del Problema (evitando la caja negra)}

Explotar la estructura del problema es común en diversas áreas de la optimización. La optimización de tipo \textit{gray-box} es un concepto relativamente nuevo que se centra en incorporar la estructura del problema en el proceso de optimización. Este enfoque busca descubrir y explorar un “orden oculto”, un concepto ampliamente estudiado en computación evolutiva en el contexto del aprendizaje de dependencias (\textit{linkage learning}). Las dependencias entre variables son cruciales para el rendimiento de los algoritmos genéticos (GAs). Sin conocer estas dependencias, incluso problemas separables pueden ser exponencialmente difíciles para GAs simples, y su resolución puede requerir tamaños de población exponencialmente grandes para encontrar el óptimo global.

En la optimización \textit{gray-box}, se asume que la estructura del problema es conocida \textit{a priori}, como ocurre en problemas discretos y combinatorios.

Sin embargo, en problemas continuos, la estructura puede no ser evidente, por lo que primero debe ser descubierta. Se han propuesto diversos algoritmos para analizar la interacción entre variables y capturar la topología del problema, lo cual no solo convierte problemas de tipo \textit{black-box} en \textit{gray-box}, sino que también permite revelar información no trivial, incluso en problemas completamente conocidos (\textit{white-box}).

La información estructural puede emplearse de diferentes formas para mejorar el rendimiento en la optimización. Algunos métodos, como la \textit{cooperative co-evolution} (CC) \cite{CC}, descomponen el problema en subproblemas de menor dimensionalidad, requiriendo un análisis de interacción de variables para lograr una descomposición adecuada. Otros métodos, como los algoritmos evolutivos basados en distribuciones (\textit{EDAs}) y los algoritmos de optimización bayesiana (\textit{BOAs}), no descomponen el problema explícitamente, pero capturan y aprovechan la información de interacción mediante modelos probabilísticos construidos durante el proceso de optimización.


Para descubrir la estructura interna del problema, existen dos aproximaciones, los métodos implícitos y los métodos explícitos. A continuación se presenta una clasificación de ambos métodos:

\begin{table}[H]
\renewcommand{\arraystretch}{1.5} % Ajusta la altura de las filas
\centering
\resizebox{\textwidth}{!}{%
\begin{tabular}{|l|p{12cm}|}
\hline
\textbf{Método} & \textbf{Descripción} \\ \hline

\textbf{Adaptación de Interacciones} & Métodos que extienden algoritmos genéticos (GAs) simples agregando mecanismos para mejorar la representación del problema y promover una vinculación estrecha (\textit{tight linkage}). \\ \hline

\textbf{Modelos Probabilísticos} & Métodos que utilizan distribuciones de probabilidad para representar la función objetivo o características del problema. \\ \hline
\multicolumn{2}{|l|}{\textbf{Submétodos:}} \\ \hline

- \textit{Construcción de un modelo de interacción de variables} & Los \textit{EDAs} y \textit{BOAs} modelan las interacciones entre variables mediante redes bayesianas o matrices de covarianza adaptativas. Este enfoque captura patrones complejos en el problema. \\ \hline

- \textit{Construcción de un modelo de la función objetivo} & La optimización bayesiana construye un modelo probabilístico de la función objetivo basado en distribuciones previas y posteriores, que se actualizan con nuevas evaluaciones. \\ \hline

- \textit{Construcción de un modelo del movimiento poblacional} & Estrategias como \textit{CMA-ES} modelan el movimiento poblacional utilizando distribuciones gaussianas multivariantes, ajustando la matriz de covarianza para reflejar el paisaje del problema. \\ \hline

\textbf{Reducción Dimensional y Partición del Espacio} & Técnicas que reducen la dimensionalidad o dividen el espacio en subespacios más pequeños para simplificar la complejidad computacional en problemas de gran escala. Ejemplos incluyen proyecciones aleatorias y análisis de componentes principales (PCA). \\ \hline

\textbf{Distribuciones de Cola Pesada} & Uso de distribuciones como Lévy, Cauchy o t-distributions para mejorar la exploración y la diversidad poblacional. Estas distribuciones son útiles para evitar el estancamiento en óptimos locales. \\ \hline

\end{tabular}%
}
\caption{Resumen de métodos implícitos para la explotación de la estructura del problema en optimización.}
\label{tab:implicit_methods}
\end{table}

\begin{table}[H]
\renewcommand{\arraystretch}{1.5} % Ajusta la altura de las filas
\centering
\resizebox{\textwidth}{!}{%
\begin{tabular}{|l|p{12cm}|}
\hline
\textbf{Método} & \textbf{Descripción} \\ \hline

\textbf{Descomposición en Coevolución Cooperativa (CC)} & Divide el problema en subproblemas de menor dimensionalidad, optimizando cada uno por separado en un esquema de coevolución. \\ \hline
\multicolumn{2}{|l|}{\textbf{Submétodos:}} \\ \hline
- \textit{CCGA (Cooperative Coevolution Genetic Algorithm)} & Divide un problema de $n$ dimensiones en $n$ subproblemas unidimensionales. Útil para problemas de baja interacción entre variables. \\ \hline
- \textit{Divide-in-half} & Divide el problema en dos subcomponentes de igual tamaño, optimizados iterativamente con estrategias como DE (\textit{Differential Evolution}). \\ \hline

\textbf{Gestión de Interacciones entre Variables} & Métodos que identifican relaciones entre variables para minimizar interacciones entre componentes. \\ \hline
\multicolumn{2}{|l|}{\textbf{Submétodos:}} \\ \hline
- \textit{Agrupamiento Aleatorio (Random Grouping)} & Reorganiza aleatoriamente variables después de cada ciclo evolutivo para aumentar la probabilidad de agrupar variables relacionadas. \\ \hline
- \textit{Agrupamiento Delta (Delta Grouping)} & Ordena variables según sus desplazamientos medios entre iteraciones, agrupándolas por magnitud similar. \\ \hline
- \textit{Minimización de Diferencias de Fitness (DI)} & Reorganiza variables para minimizar diferencias en las interacciones detectadas, optimizando componentes uniformes. \\ \hline

\textbf{Métodos Estadísticos} & Inferencia de interacciones utilizando análisis estadístico de la población en evolución. \\ \hline

\textbf{Métodos Basados en Diferencias Finitas} & Utilizan diferencias finitas para detectar interacciones entre pares de variables. \\ \hline
\multicolumn{2}{|l|}{\textbf{Submétodos:}} \\ \hline
- \textit{DG (Differential Grouping)} & Forma componentes no separables mediante un análisis iterativo de interacciones con menor coste computacional. \\ \hline

\textbf{Agrupamiento Automático y Semiautomático} & Algoritmos que forman grupos de forma automática o requieren que el usuario especifique el número/tamaño de los componentes. \\ \hline
\multicolumn{2}{|l|}{\textbf{Submétodos:}} \\ \hline
- \textit{Automático} & Métodos como \textit{DG2} usan matrices de interacción y algoritmos de componentes conectados para agrupar. \\ \hline
- \textit{Semiautomático} & Requieren información adicional, como el número o tamaño de componentes, para formar grupos. \\ \hline
- \textit{k-s Dimensional Components} & Métodos que requieren que el usuario especifique tanto el número como el tamaño de los componentes para formar los grupos. \\ \hline
\end{tabular}%
}
\caption{Resumen de métodos explícitos para la explotación de la estructura del problema en optimización.}
\label{tab:explicit_methods}
\end{table}

Tras haber establecido una clasificación, podemos observar que DG2 y ERDG pertenecen a la categoría de algoritmos explícitos, basados en diferencias finitas y automáticos.

\subsection{Hibridación y Algoritmos Meméticos}

El teorema de No Free Lunch \cite{no_free_lunch} indica que ningún algoritmo de búsqueda puede superar consistentemente a todos los demás en todos los problemas posibles. La \textbf{hibridación} busca combinar las fortalezas de diferentes algoritmos para mejorar su rendimiento, la calidad de las soluciones y su integración en sistemas más amplios.

\begin{itemize}
\item{Algoritmos de Búsqueda Local Híbrida}
Estos algoritmos se basan únicamente en búsqueda local sin un componente global explícito. Por ejemplo, el \textit{Multiple Trajectory Search} (MTS) emplea tres métodos de búsqueda local, seleccionando el mejor para optimizar en vecindades específicas. Este enfoque ha demostrado eficacia en problemas de hasta 1000 dimensiones.

\item{Algoritmos Meméticos}
Integran búsqueda local dentro de un marco evolutivo global, balanceando exploración y explotación. Son populares en la optimización de gran escala, destacándose en competiciones de LSGO. Un ejemplo es el algoritmo SHADE-ILS, que utilizaremos en nuestra propuesta. Los principales aspectos de diseño incluyen:
\begin{itemize}
    \item \textbf{Frecuencia de búsqueda local:} Puede ser fija, adaptativa o probabilística.
    \item \textbf{Selección de soluciones:} Basada en rendimiento, aleatoriedad o todas las soluciones.
    \item \textbf{Intensidad de búsqueda:} Determina la duración de la búsqueda local, con enfoques fijos o adaptativos.
    \item \textbf{Procedimientos de búsqueda local:} Amplia variedad, incluyendo métodos como MTS-LS, L-BFGS-B.
\end{itemize}

\item{Hibridación Coevolutiva}
Combina la descomposición de problemas con algoritmos meméticos. Los problemas se dividen en subproblemas optimizados con algoritmos globales seguidos de episodios de búsqueda local. 
\end{itemize}

\subsection{Operadores de Muestreo y Variación en DE y PSO}

Los operadores de muestreo y variación buscan mantener la diversidad en la población y mejorar la eficacia de los algoritmos en la exploración de grandes espacios de búsqueda. Dos enfoques comunes son:

\begin{itemize}

\item{Evolución Diferencial (DE)}

La Evolución Diferencial (DE) \cite{DE} es un algoritmo popular en la optimización global debido a su simplicidad y efectividad. Variantes como \textit{SHADE} y \textit{SHADE-ILS} han surgido como adaptaciones de DE para problemas de gran escala:
\begin{itemize}
    \item \textbf{SHADE}: Una variante de DE que ajusta adaptativamente el tamaño de la población y los parámetros de mutación para mantener la diversidad en poblaciones grandes.
    \item \textbf{SHADE-ILS}: Extiende SHADE mediante la integración de estrategias de búsqueda local, mejorando la precisión en problemas de alta dimensionalidad.
\end{itemize}

\item{Particle Swarm Optimization (PSO)}

PSO \cite{PSO} es un método basado en el comportamiento social de partículas. Aunque efectivo en problemas de baja dimensionalidad, enfrenta retos en alta dimensionalidad, para lo cual se han introducido estrategias de \textit{mantenimiento de diversidad} y \textit{re-inicialización}.

\end{itemize}

\begin{table}[H]
\renewcommand{\arraystretch}{1.5} % Ajusta la altura de las filas
\centering
\resizebox{\textwidth}{!}{%
\begin{tabular}{|l|p{12cm}|}
\hline
\textbf{Método} & \textbf{Descripción} \\ \hline

\textbf{Estrategias de mutación en DE} & Variaciones de la estrategia de mutación para mejorar la convergencia y diversidad en problemas de gran escala. \\ \hline
\multicolumn{2}{|l|}{\textbf{Submétodos:}} \\ \hline
- \textit{Adaptación de estrategias de mutación} & Aplicación adaptativa de estrategias según el tipo de problema, como DE/rand/1 o DE/current-to-best/1. \\ \hline
- \textit{Selección de vectores} & Uso de vectores basados en calidad o cercanía, como la combinación de global-best y personal-best. \\ \hline

\textbf{Adaptación de parámetros en DE} & Ajuste dinámico de parámetros como el factor de escala ($F$) y la tasa de cruce ($CR$) para mejorar la exploración y explotación. \\ \hline
\multicolumn{2}{|l|}{\textbf{Submétodos:}} \\ \hline
- \textit{Muestreo probabilístico} & Uso de distribuciones (uniforme, gaussiana, Cauchy) para generar parámetros adaptativos. \\ \hline
- \textit{Procesos caóticos} & Ajuste de $F$ y $CR$ mediante procesos caóticos para mejorar la búsqueda en espacios complejos. \\ \hline

\textbf{Mantenimiento de diversidad en DE} & Prevención de pérdida de diversidad en alta dimensionalidad mediante partición del espacio, coevolución o uso de archivos de soluciones. \\ \hline
\multicolumn{2}{|l|}{\textbf{Submétodos:}} \\ \hline
- \textit{Multipoblación} & Subpoblaciones con estrategias independientes para promover exploración y explotación. \\ \hline
- \textit{Archivos externos} & Almacenamiento de soluciones descartadas o fallidas para diversificar el cruce y la mutación. \\ \hline

\textbf{Actualización de PSO} & Modificación de reglas de actualización para evitar convergencia prematura y mejorar la exploración en alta dimensionalidad. \\ \hline
\multicolumn{2}{|l|}{\textbf{Submétodos:}} \\ \hline
- \textit{Reglas de actualización alternativas} & Reducción de la dependencia en global-best mediante aprendizajes sociales o mutaciones distribuidas. \\ \hline
- \textit{Topologías dinámicas} & Uso de vecindarios dinámicos o estructuras de multipoblación para mejorar la diversidad. \\ \hline

\textbf{Mantenimiento de diversidad en PSO} & Mecanismos para equilibrar la exploración y explotación en grandes dimensiones. \\ \hline
\multicolumn{2}{|l|}{\textbf{Submétodos:}} \\ \hline
- \textit{Reinicialización parcial} & Reubicación de partículas en áreas de alta actividad para concentrar la búsqueda. \\ \hline
- \textit{Muestreo basado en oposición} & Exploración en espacios opuestos para aumentar las probabilidades de mejora. \\ \hline

\textbf{Partición del espacio en PSO} & División del espacio de búsqueda en subregiones optimizadas por separado para mejorar el enfoque y evitar estancamientos. \\ \hline
\multicolumn{2}{|l|}{\textbf{Submétodos:}} \\ \hline
- \textit{Agrupamiento de dimensiones} & Actualización segmentada de dimensiones para evitar convergencia prematura. \\ \hline
- \textit{Subenjambres} & Subenjambres independientes que comparten información de forma controlada para mantener la diversidad. \\ \hline

\end{tabular}%
}
\caption{Resumen de operadores de muestreo y variación en DE y PSO para optimización global a gran escala.}
\label{tab:sampling_variation}
\end{table}

Una vez terminada una exposición de donde se sitúan los algoritmos que vamos a utilizar, discutiremos sus propiedades teóricas y realizaremos un estudio comparativo para comprobar si el agrupamiento diferencial de variables es efectivo en todos los casos. Llegaremos a la conclusión de que no se puede aplicar directamente la descomposición de variables con cualquier algoritmo. Los resultados de SHADE-ILS, demuestran que en algoritmos pensados para al alta dimensionalidad, aplicar descomposición no necesariamente mejora los resultados obtenidos, y es necesario realizar mejoras en este algoritmo si se quiere utilizar con problemas descompuestos. Esto último, se deja como trabajo futuro.

\endinput
               

% -------------------------------------------------------------------
% MAINMATTER
% -------------------------------------------------------------------
\mainmatter % activa la numeración de capítulos, resetea la numeración de las páginas y usa números arábigos

\part{Parte Matemática} % Dividir un TFG en partes OPCIONAL

% Información relevante para la elaboración del trabajo.
% !TeX root = ../tfg.tex
% !TeX encoding = utf8

\chapter{Documentación}\label{ch:primer-capitulo}

\section{Introducción}
Este documento es una plantilla para la elaboración de un trabajo fin de Grado siguiendo los \href{https://grados.ugr.es/matematicas/pages/infoacademica/tfg/requisitosTFG}{requisitos} de la comisión de Grado en Matemáticas de la Universidad de Granada que, a fecha de junio de 2023, son las siguientes:

\begin{itemize}
  \item La  memoria  debe  realizarse  con  un  procesador  de  texto  científico,  preferiblemente (La)TeX.
  \item La portada  debe contener  el  logo  de  la UGR,  incluir  el  título del TFG, el nombre del estudiante y especificar el grado, la facultad y el curso actual.
  \item La contraportada contendrá además el nombre del tutor o tutores.
  \item La memoria debe necesariamente incluir:
    \begin{itemize}
      \item Declaración explícita firmada en la que se asume la originalidad del trabajo, entendida en el sentido de que no ha utilizado fuentes sin citarlas debidamente. Esta declaración se puede descargar en la web del Grado.
      \item un índice detallado de capítulos y secciones,
      \item un resumen amplio en inglés del trabajo realizado (se recomienda entre 800 y 1500 palabras),
      \item una introducción en la que se describan claramente los objetivos previstos inicialmente en la propuesta de TFG, indicando si han sido o no alcanzados, los antecedentes importantes para el desarrollo, los resultados obtenidos, en su caso y las principales fuentes consultadas,
      \item una bibliografía final que incluya todas las referencias utilizadas.
    \end{itemize}
  \item Se recomienda que la extensión de la memoria sea de unas 50 páginas, sin incluir posibles apéndices.
\end{itemize}

Para generar el pdf a partir de la plantilla basta compilar el fichero \texttt{tfg.tex}. Es conveniente leer los comentarios contenidos en dicho fichero pues ayudarán a entender mejor como funciona la plantilla. 

La estructura de la plantilla es la siguiente\footnote{Los nombres de las carpetas no se han acentuado para evitar problemas en sistemas con Windows}: 
\begin{itemize}
  \item Carpeta \textbf{preliminares}: contiene los siguientes archivos
    \begin{description}
      \item[\texttt{dedicatoria.tex}] Para la dedicatoria del trabajo (opcional)
      \item[\texttt{agradecimientos.tex}] Para los agradecimientos del trabajo (opcional)
      \item[\texttt{introduccion.tex}] Para la introducción (obligatorio)
      \item[\texttt{summary.tex}] Para el resumen en inglés (obligatorio)
      \item[\texttt{tablacontenidos.tex}] Genera de forma automática la tabla de contenidos, el índice de figuras y el índice de tablas. Si bien la tabla de contenidos es conveniente incluirla, el índice de figuras y tablas es opcional. Por defecto está desactivado. Para mostrar dichos índices hay que editar este fichero y quitar el comentario a \verb+\listoffigures+ o \verb+\listoftables+ según queramos uno de los índices o los dos. En este archivo también es posible habilitar la inclusión de un índice de listados de código (si estos han sido incluidos con el paquete \texttt{listings})
  \end{description}
  El resto de archivos de dicha carpeta no es necesario editarlos pues su contenido se generará automáticamente a partir de los metadatos que agreguemos en \texttt{tfg.tex}

  \item Carpeta \textbf{capitulos}: contiene los archivos de los capítulos del TFG. Añadir tantos archivos como sean necesarios. Este capítulo es \texttt{capitulo01.tex}.

  \item Carpeta \textbf{apendices}: Para los apéndices (opcional)
  \item Carpeta \textbf{img}: Para incluir los ficheros de imagen que se usarán en el documento.
    
  \item Fichero \texttt{library.bib}: Para incluir las referencias bibliográficas en formato \texttt{bibtex}. Es útil la herramienta \href{https://www.doi2bib.org/}{doi2bib} para generar de forma automática el código bibtex de una referencia a partir de su \textsc{doi}  así como la base de datos bibliográfica \href{https://mathscinet.ams.org}{MathSciNet}. Para que una referencia aparezca en el pdf no basta con incluirla en el fichero \texttt{library.bib}, es necesario además \emph{citarla} en el documento usando el comando \verb+\cite+. Si queremos mostrar todos las referencias incluidas en el fichero \texttt{library.bib} podemos usar \verb+\cite{*}+ aunque esta opción no es la más adecuada. Se aconseja que los elementos de la bibliografía estén citados al menos una vez en el documento (y de esa forma aparecerán de forma automática en la lista de referencias).

  \item Fichero \texttt{glosario.tex}: Para incluir un glosario en el trabajo (opcional). Si no queremos incluir un glosario deberemos borrar el comando \verb+% !TeX root = ../tfg.tex
% !TeX encoding = utf8

\chapter*{Glosario}
\addcontentsline{toc}{chapter}{Glosario} % Añade el glosario a la tabla de contenidos

La inclusión de un glosario es opcional.

Archivo: \texttt{glosario.tex}

\begin{description} 
  \item[$\mathbb{R}$] Conjunto de números reales.

  \item[$\mathbb{C}$] Conjunto de números complejos.

  \item[$\mathbb{Z}$] Conjunto de números enteros.
\end{description}
\endinput
+ del fichero \texttt{tfg.tex} y posteriormente borrar el fichero \texttt{glosario.tex}

   \item Fichero \texttt{tfg.tex}: El documento maestro del TFG que hay que compilar con \LaTeX\ para obtener el pdf. En dicho documento hay que cambiar la \emph{información del título del \textsc{tfg} y el autor así como los tutores}.
\end{itemize}



\section{Elementos del texto}

En esta sección presentaremos diferentes ejemplos de los elementos de texto básico. Conviene consultar el contenido de \texttt{capitulos/capitulo01.tex} para ver cómo se han incluido.

\subsection{Listas}
En \LaTeX\ tenemos disponibles los siguientes tipos de listas:

Listas enumeradas:
\begin{enumerate}
  \item item 1
  \item item 2
  \item item 3
\end{enumerate}

Listas no enumeradas
\begin{itemize}
  \item item 1
  \item item 2
  \item item 3
  \end{itemize}

Listas descriptivas
\begin{description}
  \item[termino1] descripción 1
  \item[termino2] descripción 2
\end{description}
  
\subsection{Tablas y figuras}

En la \autoref{tb:ejemplo-tabla} o la \autoref{fig:logo-ugr} podemos ver\ldots

\begin{table}[htpb]
  \centering
  \begin{tabular}{ccc} \toprule
    \multicolumn{2}{c}{Agrupados} \\ \cmidrule(r){1-2}
    cabecera & cabecera & cabecera          \\ \midrule
    elemento & elemento & elemento          \\ 
    elemento & elemento & elemento          \\ 
    elemento & elemento & elemento          \\ \bottomrule
  \end{tabular}
  \caption{Ejemplo de tabla}
  \label{tb:ejemplo-tabla}
\end{table}

\begin{figure}[htpb]
  \centering
  \includegraphics[width=0.5\textwidth]{logo-ugr}
  \caption{Logotipo de la Universidad de Granada}
  \label{fig:logo-ugr}
\end{figure}

\section{Entornos matemáticos}\label{sec:entornos-matematicos}

La plantilla tiene definidos varios entornos matemáticos cuyo nombre es el mismo omitiendo los acentos. Así, para incluir una \emph{proposición} usaríamos:

\begin{verbatim}
\begin{proposicion}
texto de la proposición
\end{proposicion} 
\end{verbatim}

Ver el código fuente del archivo \texttt{documentacion.tex} en la carpeta \texttt{capitulos} para el resto de ejemplos.

\begin{teorema}\label{thm:teorema}
Esto es un ejemplo de teorema.
\end{teorema}

\begin{proposicion}
Ejemplo de proposición
\end{proposicion}

\begin{lema}
Ejemplo de lema
\end{lema}

\begin{corolario}
Ejemplo de corolario
\end{corolario}

\begin{definicion}
Ejemplo de definición
\end{definicion}

\begin{observacion}
Ejemplo de observación
\end{observacion}

Adicionalmente está definido el entorno \texttt{teorema*} que permite incluir un teorema sin numeración:

\begin{teorema*}[Fórmula de Gauß-Bonnet]
  Sea $S$ una superficie compacta y $K$ su curvatura de Gauß. Entonces
\begin{equation}
  \int_S K = 2\pi\chi(S)
\end{equation}
donde $\chi(S)$ es la característica de Euler de $S$.
\end{teorema*}

Las fórmulas matemáticas se escriben entre símbolos de dólar \$ si van en línea con el texto o bien usando el entorno%
\footnote{
  También es posible delimitar una ecuación mediante los comandos \texttt{$\backslash$[} y \texttt{$\backslash$]} pero éstas nunca llevarán numeración aunque añadamos una etiqueta y las referenciemos (ver \autoref{sec:referencias}).
} 
\texttt{equation} cuando queremos que se impriman centradas en una línea propia, como el siguiente ejemplo
\begin{equation}\label{eq:identidad-pitagorica}
  \cos^2 x + \sin^2 x = 1.
\end{equation}


Gracias al paquete \texttt{mathtools}, las ecuaciones escritas dentro del entorno \texttt{equation} llevarán numeración de forma automática si son referenciadas  en cualquier parte del documento (por ejemplo la identidad Pitagórica~\eqref{eq:identidad-pitagorica}, ver el código de los dos anteriores ejemplos y la \autoref{sec:referencias} para más información sobre referencias cruzadas en el documento).

\section{Listados de código}

Podemos incluir un archivo externo de código mediante el comando \texttt{lstinputlisting} especificando su nombre completo (incluyendo la extensión) y usando la opción \texttt{inputpath} para indicar la ruta hacia el fichero (siempre referida a la carpeta principal de la plantilla) así como la opción \texttt{language} para indicar el lenguaje de programación en que está escrito (esto permitirá a \LaTeX\ colorear adecuadamente el código). Además, si lo consideramos necesario, podemos indicar las líneas que queremos mostrar (ver el código fuente del \autoref{code:prime}). Consultar todas las opciones posibles en la \href{https://osl.ugr.es/CTAN/macros/latex/contrib/listings/listings.pdf}{documentación del paquete \texttt{listings}}.

\lstinputlisting[inputpath=code, language=R, linerange={11-17}, firstnumber={11}, caption={Extracto código (líneas de 11 a 17) del fichero \texttt{primeR.r}}, label={code:prime}]{primeR.r}

Alternativamente, podemos incluir el código en un entorno \texttt{lstlisting} como el \autoref{code:perceptron}

\begin{lstlisting}[caption={Implementación de un perceptrón}, label={code:perceptron}, language={python}]
def dot(v, w):
    """Producto escalar de v y w, |$\color{comment}v_0 \cdot   w_0 + \cdots + v_n \cdot w_n$|"""
    return sum(v_i * w_i for v_i, w_i in zip(v, w))

def funcion_activacion(x):
    """1 si la entrada es mayor o igual que 1, 0 en otro caso."""
    return 1 if x >= 0 else 0

def perceptron(entrada, pesos):
    """1 si el perceptron se activa, 0 en otro caso"""
    return funcion_activacion(dot(entrada, pesos))
\end{lstlisting}

La opción \texttt{float} al incluir un listado de código permitará a dicho bloque ``flotar'' como si fuese un entorno \texttt{figure} y de esta manera evitaremos que se corte al final de una página.



\section{Referencias a elementos del texto}\label{sec:referencias}

Para las referencias a los elementos del texto (secciones, capítulos, teoremas,\ldots) se procede de la siguiente manera:
\begin{itemize}
  \item Se \emph{marca} el elemento (justo después del mismo si se trata de un capítulo o sección o en el interior del \emph{entorno} en otro caso), mediante el comando \verb+\label{+\emph{etiqueta}\verb+}+, donde \emph{etiqueta} debe ser un identificador único.
  \item Para crear una referencia al elemento en cualquier otra parte del texto se usa el comando \verb+\ref{+\emph{etiqueta}\verb+}+ (únicamente imprime la numeración asociada a dicho elemento, por ejemplo \ref{ch:primer-capitulo} o \ref{thm:teorema}) o bien \verb+\autoref{+\emph{etiqueta}\verb+}+ (imprime la numeración del elemento así como un texto previo indicando su tipo, por ejemplo \autoref{ch:primer-capitulo} o \autoref{thm:teorema})
\end{itemize}




\section{Bibliografía e índice}

Esto es un ejemplo de texto en un capítulo. Incluye varias citas tanto a libros~\cite{Aigner2018}, artículos de investigación~\cite{Euler1985}, recursos online~\cite{EulerWiki} (páginas web), tesis~\cite{CitekeyPhdthesis}, trabajo fin de máster~\cite{CitekeyMastersthesis}, trabajo fin de grado~\cite{CiteKeyBachelorsthesis} así como artículos sin publicar (preprints) \cite{castroinfantes2022conjugate} (en estos últimos usar el campo \texttt{note} para añadir la información relevante). 

Ver el fichero \texttt{library.bib} para las distintas plantillas. Cada nueva referencia debe añadirse en dicho fichero siguiendo el estilo del código \texttt{bibtex} según el tipo de referencia (página web, tesis, trabajo fin de grado o máster, artículo de investigación, libro,\ldots). Alternativamente se puede usar la web \href{https://zbib.org}{https://zbib.org} para generar automáticamente el código \texttt{bibtex}.


\endinput

\chapter[Consideraciones elaboración TFG]{Consideraciones generales para la elaboración de un trabajo fin de grado}

\section{Normativa de la comisión del Grado en Matemáticas}

El \textsc{tfg} lo rigen dos normativas: 
\begin{itemize}
  \item una a nivel general de la UGR (\href{https://secretariageneral.ugr.es/sites/webugr/secretariageneral/public/inline-files/BOUGR/187/PLANTILLA%20CABECERASDoc2.pdf}{Reglamento del Trabajo o Proyecto fin de Grado de la Universidad de Granada}\footnote{\url{https://secretariageneral.ugr.es/sites/webugr/secretariageneral/public/inline-files/BOUGR/187/PLANTILLA\%20CABECERASDoc2.pdf}}) y 
  \item otra complementaria a nivel de la Facultad de Ciencias  (\href{https://fciencias.ugr.es/images/stories/documentos/reglamentos/reglamentoTfgCiencias23.pdf}{Reglamento del trabajo fin de grado en la Facultad de Ciencias de la Universidad de Granada}\footnote{\url{https://fciencias.ugr.es/images/stories/documentos/reglamentos/reglamentoTfgCiencias23.pdf}}). 
\end{itemize} 
Además, la comisión del Grado de Matemáticas impone unos \href{https://grados.ugr.es/matematicas/pages/infoacademica/tfg/requisitosTFG}{Requisitos de la memoria}\footnote{\url{https://grados.ugr.es/matematicas/pages/infoacademica/tfg/requisitosTFG}}.

El \textsc{tfg} hay que elaborarlo preferiblemente en LaTeX y puede usar la plantilla disponible en \href{https://github.com/latex-mat-ugr/Plantilla-TFG/archive/master.zip}{Plantilla \textsc{tfg} grado en matemáticas formato .tex}\footnote{\url{https://github.com/latex-mat-ugr/Plantilla-TFG/archive/master.zip}}.

Toda la información anterior puede encontrarse en la \href{https://grados.ugr.es/matematicas/pages/infoacademica/trabajofingrado}{web del Grado en Matemáticas}\footnote{\url{https://grados.ugr.es/matematicas/pages/infoacademica/trabajofingrado}}.

Es conveniente tener presente la documentación anterior para la elaboración del \textsc{tfg}. En especial en lo relativo a las fechas de depósito del \textsc{tfg} para su defensa.

A continuación destaco algunos aspectos importantes de la misma:
\begin{itemize}
\item El plagio, entendido como la presentación de un trabajo u obra hecho por otra persona como propio o la copia de textos sin citar su procedencia y dándolos como de elaboración propia, conllevará automáticamente la calificación numérica de cero. Esta consecuencia debe entenderse sin perjuicio de las responsabilidades disciplinarias en las que pudieran incurrir los estudiantes que plagien.
  \item Las memorias entregadas por parte de los estudiantes tendrán que ir firmadas sobre una declaración explícita en la que se asume la originalidad del trabajo, entendida en el sentido de que no ha utilizado fuentes sin citarlas debidamente.
\end{itemize}


% Los criterios de evaluación utilizados permitirán evaluar el grado de adquisición de las competencias que tiene establecidas el TFG en el VERIFICA de la titulación. Además, entre otros aspectos, se tendrá en consideración:

% \begin{itemize}
%   \item Redacción y ortografía tanto en la memoria del TFG como en los medios usados para la defensa del mismo (diapositivas, etc.).
%   \item Adecuación al formato de memoria indicado. Se proporcionará una plantilla de memoria de TFG a tal fin.
%   \item Adecuación temporal a los cronogramas de trabajo según los plazos de entrega marcados por el tutor/es.
%   \item Nivel de profundidad en los contenidos expuestos.
%   \item Dominio del tema e iniciativa del alumno.
%   \item Claridad de la exposición y adecuación al tiempo de exposición establecido.
%   \item Capacidad de análisis y síntesis.
%   \item Discusión con la Comisión Evaluadora en el turno de preguntas.
% \end{itemize}

% Se entregará una copia escrita, a doble cara, de la memoria del TFG para su evaluación, por parte de la Comisión Evaluadora, con una antelación de una semana (7 días naturales) antes de la fecha de defensa pública del TFG. Además, se entregará una versión en formato "pdf" de dicha memoria que quedará en la base de datos de todos los TFG y que custodiará la CTFGOO.




\section{Formato de la memoria}
La memoria se presentará usando un editor de textos científico, preferiblemente \LaTeX, e incluir los siguientes apartados:
\begin{enumerate}
  \item \emph{Resumen en inglés}: Deberá estar escrito completamente en inglés y tener una longitud recomendada entre 800 y 1500 palabras. 
  \item \emph{Introducción}. Deberá:
    \begin{itemize}
      \item Indicar los \emph{Objetivos del trabajo}: deberán aparecer con claridad los objetivos inicialmente previstos en la propuesta de \textsc{tfg} y los finalmente alcanzados con indicación de dificultades, cambios y mejoras respecto a la propuesta inicial. Si procede, es conveniente apuntar de manera precisa las interdependencias entre los distintos objetivos y conectarlos con los diferentes apartados de la memoria. Se pueden destacar aquí los aspectos formativos previos más utilizados. 
    \item Contextualizar el trabajo explicando antecedentes importantes para el desarrollo realizado y efectuando, en su caso, un estudio de los progresos recientes.
    \item Describir el problema abordado, de forma que el lector tenga desde este momento una idea clara de la cuestión a resolver o del producto a desarrollar y una visión general de la solución alcanzada.
    \item Indicar los resultados obtenidos.
    \item Citar las principales fuentes consultadas.
    \end{itemize}

  \item \emph{Desarrollo del trabajo}: El trabajo se estructurará en partes o capítulos según convengan, con la posibilidad de incluir apéndices. Se recomienda que la extensión de esta parte (sin incluir los apéndices) sea de unas 50 páginas. 

  \item \emph{Conclusiones y vías futuras}: Las conclusiones deberán incluir todas aquellas de tipo profesional y académico. Si hubiese posibles vías claras de desarrollo posterior sería interesante destacarlas aquí, poniéndolas en valor en el contexto inicial del trabajo.

  \item \emph{Bibliografía final}: Se incluirán tanto las fuentes primarias como todas aquellas cuyo peso haya sido menor en la realización del trabajo. Se recomienda un breve comentario de las referencias, ya sea individualizado, por grupos de referencias o global. En caso de incluir \textsc{url}s de páginas web deberán ir acompañadas de título, autor y fecha de último acceso, entre otros datos relevantes. Se recomienda no abusar de este tipo de fuentes.
\end{enumerate}




\section{Recomendaciones}
A la hora de abordar un trabajo como este, de cierta complejidad y extensión, es conveniente tener ciertas consideraciones desde un principio que ayuden a la organización y realización del mismo.
\begin{itemize}
    \item La memoria deberá ceñirse a las directrices dadas en la sección precedente. 
    
    \item Cualquier consulta externa (libro, artículo, página web, imagen,\ldots) debe estar debidamente referenciada tanto en el texto como en la bibliografía al final del trabajo. La bibliografía debe de aparecer en orden alfabético (del primer autor) en el formato indicado en la plantilla. 

    \item Se debe evitar copiar texto de forma literal, salvo citas literales, que se indicarán como tales y entrecomilladas. LaTeX proporciona el entorno \texttt{quote} para ello.

    \item Todas las imágenes y tablas incluidas en el documento deben figurar con su respectivos créditos (excepto que sean de elaboración propia). Por tanto, es recomendable guardar las referencias consultadas (direcciones web, libros) para la obtención de cualquier material gráfico o de datos.

    \item Si el trabajo contiene gran cantidad de vocabulario específico, conviene añadir un glosario de términos al final del mismo. Esto es mejor ir haciéndolo conforme se avanza en la redacción del trabajo.

    \item Es conveniente hacer un esquema inicial con la estructura general de la memoria: ¿de cuántas partes constará? ¿en qué orden? ¿qué incluirá cada una de ellas? En la plantilla proporcionada se recomienda una estructura general. Ello ayudará a organizar mejor el trabajo. No obstante, dicha estructura inicial puede ser modificada cuando el trabajo esté avanzado si el contenido lo requiere.
\end{itemize}


% Añadir tantos capítulos como sea necesario

\cleardoublepage\part{Parte Informática}
% !TeX root = ../tfg.tex
% !TeX encoding = utf8

\chapter{Ejemplo de capítulo}

\section{Primera sección}

Este fichero \texttt{capitulo-ejemplo.tex} es una plantilla para añadir capítulos al \textsc{tfg}. Para ello, es necesario:
\begin{itemize}
  \item Crear una copia de este fichero \texttt{capitulo-ejemplo.tex} en la carpeta \texttt{capitulos} con un nombre apropiado (p.e. \texttt{capitulo01.tex}).
  \item Añadir el comando \texttt{$\backslash$input\{capitulos/capitulo01\}} en el fichero principal \texttt{tfg.tex} donde queremos que aparezca dicho capítulo.
\end{itemize}


\endinput
%--------------------------------------------------------------------
% FIN DEL CAPÍTULO. 
%--------------------------------------------------------------------


% -------------------------------------------------------------------
% APPENDIX: Opcional
% -------------------------------------------------------------------

\appendix % Reinicia la numeración de los capítulos y usa letras para numerarlos
\pdfbookmark[-1]{Apéndices}{appendix} % Alternativamente podemos agrupar los apéndices con un nuevo \part{Apéndices}

% !TeX root = ../tfg.tex
% !TeX encoding = utf8

\chapter{Ejemplo de apéndice}\label{ap:apendice1}

Los apéndices son opcionales.

Este fichero \texttt{apendice-ejemplo.tex} es una plantilla para añadir apéndices al \textsc{tfg}. Para ello, es necesario:
\begin{itemize}
  \item Crear una copia de este fichero \texttt{apendice-ejemplo.tex} en la carpeta \texttt{apendices} con un nombre apropiado (p.e. \texttt{apendice01.tex}).
  \item Añadir el comando \texttt{$\backslash$input\{apendices/apendice01\}} en el fichero principal \texttt{tfg.tex} donde queremos que aparezca dicho apéndice (debe de ser después del comando \texttt{$\backslash$appendix}).
\end{itemize}

\endinput
%------------------------------------------------------------------------------------
% FIN DEL APÉNDICE. 
%------------------------------------------------------------------------------------

% Añadir tantos apéndices como sea necesario 

% -------------------------------------------------------------------
% GLOSARIO: Opcional
% -------------------------------------------------------------------

% !TeX root = ../tfg.tex
% !TeX encoding = utf8

\chapter*{Glosario}
\addcontentsline{toc}{chapter}{Glosario} % Añade el glosario a la tabla de contenidos

La inclusión de un glosario es opcional.

Archivo: \texttt{glosario.tex}

\begin{description} 
  \item[$\mathbb{R}$] Conjunto de números reales.

  \item[$\mathbb{C}$] Conjunto de números complejos.

  \item[$\mathbb{Z}$] Conjunto de números enteros.
\end{description}
\endinput
 

% -------------------------------------------------------------------
% BACKMATTER
% -------------------------------------------------------------------

\backmatter % Desactiva la numeración de los capítulos
\pdfbookmark[-1]{Referencias}{BM-Referencias}

% BIBLIOGRAFÍA
%-------------------------------------------------------------------

\bibliographystyle{alpha-es} 
\begin{small} 
  \bibliography{library.bib}
\end{small}


\end{document}
